\chapter{Conclusion}

In conclusion, Wayclip is an emerging ecosystem of tools, aiming to deliver a competitive experience for clipping your gameplay on Linux.
Linux, or in particular Wayland, presents a challenging environment in which to build a tool that upholds modern performance standards.
Wayclip aims to provide the better of both worlds -- performance, and ease of use. This will be achieved by using Rust, alongside powerful libraries such as GStreamer and PipeWire.

As of October 2026, Wayclip is still under development with a lot of plans ahead.
Table~\ref{tab:project_status} displays the status and publicity of each component in Wayclip.

\begin{table}[htbp]
	\centering
	\begin{tabular}{lll}
		\hline
		\textbf{Project} & \textbf{Status}               & \textbf{Publicity}               \\ \hline
		Wayclip Daemon   & Active (Open Testing)         & Public / Open-Source             \\
		Wayclip Core     & Active (Open Testing)         & Public / Open-Source             \\
		Wayclip CLI      & Testing (Closed Testing)      & Unreleased / Planned Open-Source \\
		Wayclip App      & In Development (Tauri v3 CEF) & Unreleased / Planned Open-Source \\
		Wayclip Website  & In Development (SvelteKit)    & Private / Protected (undecided)  \\
		Wayclip API      & In Development (Actix-web)    & Private / Protected (undecided)  \\ \hline
	\end{tabular}
	\caption{Development status and source publicity across Wayclip ecosystem components.}
	\label{tab:project_status}
\end{table}

The future of Wayclip seems bright, with a lot of room to expand, grow and mature if the right approach is taken in the early stages of development.

\bigskip

With that said, this paper comes to end.
